% config/listings.tex

% --- Palette Pastello "Cyber-Nerd" ---
\definecolor{nerd_bg}{RGB}{245, 245, 250}
\definecolor{nerd_keyword}{RGB}{135, 95, 215}
\definecolor{nerd_comment}{RGB}{120, 150, 120}
\definecolor{nerd_string}{RGB}{215, 95, 135}
\definecolor{nerd_num}{RGB}{95, 135, 175}
\definecolor{nerd_frame}{RGB}{200, 210, 230}
\definecolor{nerd_line}{RGB}{150, 160, 190}

\usepackage{inconsolata}

\lstset{
    language=C++,
    % --- Fix UTF-8 for "Invalid byte sequence" ---
    extendedchars=true,
    inputencoding=utf8,
    % Mappatura manuale dei caratteri che causano il crash (0xB0 è il grado, 0xC0 la À)
    literate=
        {à}{{\`a}}1 {è}{{\'e}}1 {é}{{\'e}}1 {ì}{{\`i}}1 {ò}{{\`o}}1 {ù}{{\`u}}1
        {À}{{\`A}}1 {È}{{\`E}}1 {É}{{\'E}}1 {Ì}{{\`I}}1 {Ò}{{\`O}}1 {Ù}{{\`U}}1
        {°}{{$^\circ$}}1 { }{{~}}1
        {->}{{$\rightarrow$}}2 {>=}{{$\ge$}}2 {<=}{{$\le$}}2 {==}{{$==$}}2 {!=}{{$\neq$}}2,
    % ---------------------------------------------
    backgroundcolor={},   
    basicstyle=\fontfamily{zi4}\selectfont\footnotesize\linespread{1.05},
    commentstyle=\color{nerd_comment}\itshape,
    keywordstyle=\color{nerd_keyword}\bfseries,
    numberstyle=\tiny\color{nerd_line},
    stringstyle=\color{nerd_string},
    identifierstyle=\color{black!80},
    numbers=none,                     
    numbersep=12pt,                   
    xleftmargin=30pt,                 
    framexleftmargin=5pt,
    tabsize=2,
    breaklines=true,
    breakatwhitespace=false,
    frame=single, 
    framerule=1pt,
    rulecolor=\color{nerd_frame},
    captionpos=b,
    abovecaptionskip=12pt,
    emph={setup, loop, PIN_LED, Serial, LittleFS, GNSS, uint16_t, uint32_t, uint, alarm_id_t, pio_sm_put_blocking},
    emphstyle=\color{nerd_num}\bfseries
}

\newcommand{\nt}[1]{\texttt{\color{nerd_keyword}#1}}




\definecolor{codebg}{rgb}{0.96,0.96,0.96}
\definecolor{codekey}{rgb}{0.0,0.33,0.62}
\definecolor{codecomment}{rgb}{0.28,0.50,0.28}
\definecolor{codestring}{rgb}{0.58,0.10,0.10}

\lstset{
	language=C++,
	backgroundcolor=\color{codebg},
	basicstyle=\ttfamily\footnotesize,
	keywordstyle=\color{codekey}\bfseries,
	commentstyle=\color{codecomment}\itshape,
	stringstyle=\color{codestring},
	breaklines=true,
	frame=single,
	rulecolor=\color{gray!35},
	numbers=left,
	numberstyle=\tiny\color{gray},
	tabsize=4,
	showstringspaces=false,
	morekeywords={uint8_t,uint32_t,int8_t,bool,size_t,int_least32_t,
		spin_lock_t,SystemDataPacket,WeatherDataPacket,
		error_category_t,error_code_t,error_report_t,
		GpsStatus_t,SystemConfig,view_id_t,view_interface_t,
		button_t,button_state_t,ButtonId,LedAnimState,
		FeedbackType,power_mode_t,TaskCallback,Task,Vector3D_t,
		CRGB,PROGMEM}
}

\lstdefinestyle{makefile}{
	language=make,
	backgroundcolor=\color{codebg},
	basicstyle=\ttfamily\footnotesize,
	commentstyle=\color{codecomment}\itshape,
	breaklines=true, frame=single,
	rulecolor=\color{gray!35}, numbers=none
}

\lstset{
	backgroundcolor=\color{codebg},
	basicstyle=\ttfamily\footnotesize,
	keywordstyle=\color{codekey}\bfseries,
	commentstyle=\color{codecomment}\itshape,
	stringstyle=\color{codestring},
	breaklines=true, frame=single,
	rulecolor=\color{gray!35},
	numbers=left, numberstyle=\tiny\color{gray},
	tabsize=4, showstringspaces=false
}
\lstdefinestyle{ini}{
	language={}, numbers=none,
	backgroundcolor=\color{codebg},
	basicstyle=\ttfamily\footnotesize,
	commentstyle=\color{codecomment}\itshape,
	stringstyle=\color{codestring},
	breaklines=true, frame=single,
	rulecolor=\color{gray!35}
}
\lstdefinestyle{py}{
	language=Python,
	backgroundcolor=\color{codebg},
	basicstyle=\ttfamily\footnotesize,
	keywordstyle=\color{codekey}\bfseries,
	commentstyle=\color{codecomment}\itshape,
	stringstyle=\color{codestring},
	breaklines=true, frame=single,
	rulecolor=\color{gray!35},
	numbers=left, numberstyle=\tiny\color{gray}
}
